\chapter{Ferramentas, Contratos e Model Context Protocol}
\label{ch:mcp}

\ficcao O modelo escreve «consultei a reclamação» com a mesma facilidade com que escreve «um duende confessou». Para distinguir uma ação real de uma frase, precisamos de uma chamada executada, um resultado e um registo. Entram as ferramentas.

\section{O ciclo de uma chamada}
Imagina um formulário para consultar uma reclamação. Um campo diz que operação fazer e outro indica o número da reclamação. Os valores fornecidos à operação chamam-se \emph{argumentos}. Em \emph{function calling}, a aplicação descreve a ferramenta, o modelo pode propor os campos preenchidos e a aplicação valida e executa a operação. O resultado pode voltar ao modelo para fundamentar a resposta~\citep{functioncalling2026}. O formato exato varia por fornecedor.

Na Lavandaria, expomos \texttt{consultar\_estado(reclamacao\_id)}. O cliente autenticado vem da aplicação, fora dos argumentos gerados. Se deixássemos o modelo escolher livremente \texttt{cliente\_id}, uma frase «sou o cliente B» poderia tornar-se uma tentativa de mudança de identidade.

Um contrato mínimo define nome, finalidade, tipos, condições de erro, efeitos e limites. Prefere ferramentas estreitas, como consultar uma reclamação autorizada, a ferramentas universais, como executar qualquer SQL. «Apenas leitura» também exige cuidado: custo excessivo, dados privados e funções com efeitos continuam possíveis.

\section{A aplicação, o cliente e o servidor MCP}
O \emph{Model Context Protocol} (MCP) é um protocolo aberto para ligar aplicações com modelos a fontes de contexto e ferramentas. A aplicação anfitriã, ou \emph{host}, organiza o trabalho. O cliente MCP faz os pedidos e o servidor oferece as operações. As mensagens usam JSON-RPC, uma convenção para pedir a um programa que execute uma operação. O pedido indica a operação, os argumentos e um identificador que permite ligar a resposta ao pedido~\citep{mcp2026}.

Um servidor pode disponibilizar \emph{resources} (dados), \emph{prompts} (modelos de mensagens) e \emph{tools} (operações). São conceitos de interface. Um recurso pode conter dados pessoais; um prompt de um servidor pode conter instruções inadequadas; uma ferramenta pode ter efeitos reais. A presença num catálogo não transforma o conteúdo em fonte de confiança.

MCP não é o agente, o modelo ou a base de dados. Também não garante que diferentes clientes implementem todas as capacidades da mesma maneira. Descoberta de uma ferramenta significa saber que está exposta; autorização significa verificar se aquela identidade a pode usar naquele recurso e naquele momento.

\section{Uma versão explícita evita exemplos incompatíveis}
\begin{warningboxenv}{Referência desta edição: MCP 2026-07-28}
A especificação consultada em 3 de novembro de 2026 adota pedidos sem estado de sessão de protocolo, metadados por pedido e descoberta com \texttt{server/discover}. O changelog regista a remoção do handshake \texttt{initialize} da versão anterior~\citep{mcpchanges2026}. Exemplos de 2025 não devem ser copiados como se descrevessem esta versão.
\end{warningboxenv}

Cada pedido inclui a versão e as capacidades relevantes em \texttt{\_meta}. A ausência de estado de sessão de protocolo não impede que a aplicação guarde o estado das reclamações. Um cliente pode consultar \texttt{server/discover} para conhecer versões e capacidades suportadas; a operação é obrigatória no servidor, mas a descoberta inicial não é obrigatória em todas as chamadas. \texttt{tools/list} descreve ferramentas; \texttt{tools/call} invoca uma delas segundo o contrato da versão.

O executor da Lavandaria pode ser testado sozinho. No capítulo~\ref{ch:interop}, um segundo programa chama esse mesmo executor através de um subconjunto de MCP. A mensagem atravessa a entrada e saída de um processo Python separado, o transporte \emph{stdio}. O ensaio verifica descoberta, listagem, chamada e recusa de versão. Num produto, avalia ainda as capacidades da biblioteca de integração, ou SDK, que escolheres e testa com as implementações que terás de suportar.

\section{O limite de confiança tem de sobreviver ao prompt}
\begin{figure}[htbp]
\centering
\begin{tikzpicture}[node distance=7mm and 8mm,every node/.style={draw,align=center,text width=36mm,minimum height=10mm,font=\small},>=Stealth]
\node (model) {Modelo propõe chamada};
\node (app) [below=of model] {Aplicação valida contrato};
\node (policy) [below left=of app] {Identidade e política};
\node (exec) [below right=of app] {Executor autorizado};
\node (data) [below=25mm of app] {Dados ou efeito permitido};
\draw[->](model)--(app);\draw[->](app)--(policy);\draw[->](policy)--(exec);\draw[->](exec)--(data);
\end{tikzpicture}
\caption{A proposta do modelo não concede autorização. MCP transporta operações; a política pertence à aplicação e ao servidor.}
\label{fig:ferramentas}
\end{figure}

Em ligações remotas autenticadas, a especificação define requisitos de autorização e validação de tokens, incluindo o destinatário do token e a proibição de o passar cegamente para serviços a montante~\citep{mcpauth2026}. Segredos não pertencem ao prompt nem aos logs. Uma etiqueta de «read-only» é metadado; o executor tem de impor o comportamento.

Ao depurar uma chamada, guarda a operação pedida, a decisão de autorização, a versão do contrato e a categoria de erro. Um argumento com tipo errado deve falhar antes de consultar o recurso. Uma chamada bem formada de A para LAV-002 deve ser recusada por acesso. Um resultado de leitura pode depois falhar por falta de suporte à frase que o modelo escreveu. Registar os três casos como «erro de IA» esconde o ponto a corrigir.

Para ações com efeitos, associa a repetição a uma identidade e à mesma intenção de negócio. Se uma resposta se perder depois do efeito, confirma o estado antes de decidir se podes tentar de novo. O capítulo~\ref{ch:agentes} segue um caso com resultado desconhecido.

\begin{promptlabbox}{A ferramenta certa, com o cliente errado}
\textbf{Objetivo e contexto:} testar autorização fora do modelo. Usa o executor offline da Lavandaria: LAV-001 pertence a cliente-A e LAV-002 a cliente-B.

\textbf{Instruções:} com contexto de cliente-A, consulta LAV-001 e depois LAV-002. Tenta ainda fornecer um campo extra \texttt{cliente\_id} nos argumentos da ferramenta e usa um identificador inexistente.

\textbf{Resultado esperado:} a primeira consulta funciona; as restantes não revelam dados de outro cliente. Campo extra é rejeitado. Ausência e falta de autorização produzem uma mensagem externa equivalente.

\textbf{Validação:} executa os testes ligados a CA-02 e CA-03. Inspeciona o resultado da ferramenta, não apenas a frase final. Verifica que o cliente vem do contexto da aplicação.

\textbf{Falhas possíveis:} confiar no modelo para escolher a identidade, filtrar apenas a resposta final ou chamar a este exercício uma avaliação de compatibilidade MCP.
\end{promptlabbox}

\section{O que acrescenta uma integração real}
Além do domínio já testado, precisas de transporte, descoberta, negociação de versão, autenticação, limites, tratamento de erros e observabilidade. Valida o esquema recebido, a proveniência do servidor e o que muda quando uma ferramenta é atualizada. Testa recusa de acesso, timeout e resposta mal formada.

Se o objetivo é chamar duas funções internas, uma interface direta pode chegar. Se várias aplicações precisam de descobrir e usar as mesmas capacidades, um protocolo comum pode simplificar a integração. Escolhe-o por esse benefício, não porque três letras parecem mais inteligentes do que uma função.
